\subsection{局部凸可度量化空间的双对偶}

命题 3. --- 设 E 是一个局部凸可度量化空间，\(E_{b}^{\prime}\) 是它的强对偶，G 是一个 Fréchet 空间。空间 \(\mathcal{L}_{b}(E_{b}^{\prime}; G)\) 是 Fréchet 空间。

由（IV，第 22 页的）命题 2，存在有界子集的一个序列 \((\mathrm{A}_n)\)（\(\mathrm{E}_b'\) 中的），使得 \(\mathrm{E}_b'\) 的每个有界子集都包含在某个 \(\mathrm{A}_n\) 中。设 \((\mathrm{V}_n)\) 是 \(\mathbf{G}\) 中零邻域的一个可数基本系。设 \(\mathrm{H}_{mn}\) 是线性映射 \(u\)（从 \(\mathrm{E}_b'\) 到 \(\mathbf{G}\) 中的，且满足 \(u(\mathrm{A}_m) \subset \mathrm{V}_n\)）组成的集。则 \((\mathrm{H}_{mn})\) 是 \(\mathcal{L}_b(\mathrm{E}_b'; \mathrm{G})\) 中零邻域的一个基本系，从而后一空间是可度量化的。

为证明 \(\mathcal{L}_b(\mathrm{E}_b';\mathbf{G})\) 完备，只需证明此空间中每个 Cauchy 序列 \((u_n)\) 都收敛；由于 \(\mathbf{G}\) 完备，存在线性映射 \(u:\mathrm{E}_b' \to \mathbf{G}\) 使得 \((u_n)\) 简单收敛于 \(u\)。由 IV 第 23 页的推论，我们有 \(u \in \mathcal{L}_h(\mathrm{E}_h';\mathbf{G})\)。于是由 GT, X, § 1, 第 5 目的命题 5，\((u_n)\) 收敛于 \(u\)（在 \(\mathcal{L}_b(\mathrm{E}_b';\mathbf{G})\) 中）。

推论. --- 局部凸可度量化空间的双对偶是 Fréchet 空间。

